\chapter{对错误数据关联和离群值的鲁棒性}

\textbf{作者：} Heng Yang、Josh Mangelson、Yun Chang、Jingnan Shi、Niko Sunderhauf、Luca Carlone

在第 1 章中，我们已经看到，因子图是用于建模和可视化 SLAM 问题的强大表示；给定一组测量（例如里程计和路标测量）后，最大后验（maximum a posteriori，MAP）估计为推断感兴趣变量（例如机器人位姿和路标位置）提供了一个有理论依据且通用的框架。例如，我们看到，当测量 $z _ { i }$ 受协方差为 $\Sigma _ { i }$ 的加性零均值高斯噪声影响时，MAP 估计会导出一个非线性最小二乘优化问题：

\begin{equation}
\boldsymbol {x} ^ {\mathrm{MAP}} = \underset {\boldsymbol {x}} {\arg \min} \sum_ {i} \| \boldsymbol {z} _ {i} - \boldsymbol {h} _ {i} (\boldsymbol {x} _ {i}) \| _ {\boldsymbol {\Sigma} _ {i}} ^ {2},\tag{3.1}
\end{equation}

其中，$\pmb { x } _ { i }$ 表示参与第 $i$ 个测量的状态子集。\textsuperscript{1} 本章指出，在实践中，许多测量 $z _ { i }$ 可能由于错误的数据关联而具有很大的误差，远不符合零均值高斯分布（第 3.1 节）；这些测量通常会使式（3.1）中的估计 $\scriptstyle \mathbf { x } ^ { \mathrm { M A P } }$ 产生很大扰动。因此，我们将讨论如何在 SLAM 前端剔除严重离群值（第 3.2 节），然后重点讨论如何提高 SLAM 后端对剩余离群值的鲁棒性（第 3.3 节）。最后，我们将简要回顾近期趋势，并给出相关工作的进一步线索（第 3.4 节）。

\section{离群值的成因及其危害}

本节论证：离群值在大多数 SLAM 应用中不可避免；若未得到恰当处理，则会导致严重错误的估计。

\section{数据关联与离群值}

为了理解离群测量的成因，先考虑两个例子。首先，考虑一个基于路标的 SLAM 问题：我们需要根据里程计测量以及若干机器人位姿对路标的相对观测，重建机器人的轨迹和外部路标的位置。若如第 1 章那样假设路标测量具有零均值高斯噪声，则优化中会出现形如 $\| z _ { i j } - h ( p _ { i } , \ell _ { j } ) \| _ { \Sigma } ^ { 2 }$ 的项。这些项刻画了如下事实：给定测量 $z _ { i j }$ 是从位姿 $\pmb { p } _ { i }$ 对路标 $\ell _ { j }$ 的观测，并带有高斯噪声，其中 $h ( \cdot )$ 是描述相对测量类型（例如距离、方位等）的函数。实践中，$z _ { i j }$ 是由 SLAM 前端对原始传感器数据预处理得到的。例如，若机器人装有相机，且 $z _ { i j }$ 是对路标 $\ell _ { j }$ 方位的视觉观测，则可先在图像中进行目标（或更一般地说，特征）检测与匹配，再计算与检测到的像素对应的方位，从而提取该测量。问题在于检测并不完美：图像中被检测为 $\ell _ { j }$ 的路标，在现实中可能实际是另一个路标。这会使 $z _ { i j }$ 大幅偏离所假设的模型。将一个测量关联到某个路标的问题通常称为\textbf{数据关联问题}，它也广泛存在于其他估计问题中（例如目标跟踪）。因此，错误的数据关联会在估计问题中产生离群值。

作为第二个例子，考虑位姿图优化问题。此时我们主要关心机器人的轨迹估计（由一组位姿表示），测量则是里程计测量（关联轨迹中相邻的位姿）或回环约束（关联非相邻、且可能在时间上相距很远的位姿）。实践中，回环由基于视觉或 LiDAR 的地点识别方法检测；这些方法负责判定一对位姿 $\pmb { p } _ { i }$ 和 $p _ { j }$ 是否观测到了环境中的同一部分。遗憾的是，当前的地点识别方法容易出错，会在没有观测到同一场景的位姿之间检测出回环。这部分源于现有方法的局限性，但更多时候源于\textbf{感知混叠}（perceptual aliasing）：两个外观相似的位置实际上对应不同地点，例如大学建筑中的两个教室，或办公环境中外观相似的隔间。这同样可理解为数据关联失败：我们错误地将回环测量关联到了两个选错的机器人位姿。

需要注意，离群值不仅由错误数据关联造成，也可能由 SLAM 方法中所作假设被违背而引起。例如，大多数 SLAM 方法假设路标是静态的；因此，对运动物体的检测，即使正确地关联到该物体，也可能产生具有大残差的离群测量。类似地，传感器故障和性能退化，例如故障的轮式编码器或相机镜头上的灰尘，也可能导致测量中出现离群值。

\begin{figure}[htbp]
\centering
\includegraphics[width=\linewidth]{ddc222b702e516e02124c861c2f00219115135f9feacce05f4d6c64dcf3a872f.jpg}
\label{fig:ddc222b702e516e02124c861c2f00219115135f9feacce05f4d6c64dcf3a872f}
\end{figure}

\begin{figure}[htbp]
\centering
\includegraphics[width=\linewidth]{3ccb49c7a253eb43cabf4b89fdc57cf7a6b1778cd2652068a9406a35d9b29a0a.jpg}
\label{fig:3ccb49c7a253eb43cabf4b89fdc57cf7a6b1778cd2652068a9406a35d9b29a0a}
\end{figure}

\section{存在离群值时的最小二乘}

存在离群值时，最小二乘形式（3.1）得到的估计可能严重错误。从理论角度看，我们为测量所假设的高斯噪声是“轻尾”（light-tailed）的，即它实质上排除了误差极大的测量。从更实际的角度看，离群值会在目标函数中产生一些项：当在真值附近求值时，其残差误差 $r _ { i } ( \pmb { x } ) : = \| z _ { i j } - h ( \pmb { p } _ { i } , \pmb { \ell } _ { j } ) \| _ { \Sigma }$ 非常大。由于优化目标将残差平方，即目标为 $\textstyle \sum _ { i } r _ { i } ( { \pmb x } ) ^ { 2 }$，这些残差会对代价产生不成比例的巨大影响，优化过程会着重减小由离群值引起的大项，而非充分利用其余的内点（inlier）测量。

\begin{figure}[htbp]
\centering
\includegraphics[width=\linewidth]{dd88ae612d25550517be2aa6d8ae6d17b3901bc1d6abdd4f5d62ccc83a7c1816.jpg}
\label{fig:dd88ae612d25550517be2aa6d8ae6d17b3901bc1d6abdd4f5d62ccc83a7c1816}
\end{figure}

\begin{figure}[htbp]
\centering
\includegraphics[width=\linewidth]{14f55ee14950aebd1bd5a280ff3b53899d5efdb8ba1708bc444501ae9c30ac65.jpg}
\label{fig:14f55ee14950aebd1bd5a280ff3b53899d5efdb8ba1708bc444501ae9c30ac65}
\end{figure}

\begin{figure}[htbp]
\centering
\includegraphics[width=\linewidth]{46bebd1e4816f8ebacbb6a7544412b105d02a6f203ddb6a37bf18559bce2a3a8.jpg}
\label{fig:46bebd1e4816f8ebacbb6a7544412b105d02a6f203ddb6a37bf18559bce2a3a8}
\end{figure}

\begin{figure}[htbp]
\centering
\includegraphics[width=\linewidth]{7ffa9ac1ea7a9168e51d46d96a15361d33da886fd43be4b1d60e1058b983f8e0.jpg}
\caption{图 3.1 含离群值的 SLAM 问题：(a)–(c) 分别为 M3500、SubT 和 Victoria Park 数据集的真值轨迹；(d)–(f) 为存在离群值时采用最小二乘形式得到的轨迹估计。内点测量显示为灰色边，离群值显示为红色边。在 SubT 数据集中，还展示了由 SLAM 位姿估计构建的稠密地图。}
\label{fig:7ffa9ac1ea7a9168e51d46d96a15361d33da886fd43be4b1d60e1058b983f8e0}
\end{figure}

为说明这一点，图 3.1 给出了三个含离群值 SLAM 问题的结果。第一列是一个称为 M3500 的模拟位姿图优化基准，其中位姿按网格状构型排列；数据集包含 3500 个二维位姿和 8953 个测量。第二列是一个称为 SubT 的真实世界位姿图数据集，采集于 DARPA 地下挑战赛期间的一条隧道中 \textsuperscript{[299]}；该数据集包含 682 个三维位姿和 3278 个测量。第三列是一个称为 Victoria Park 的真实世界基于路标的 SLAM 数据集 \textsuperscript{[809]}；该数据集包含 7120 个二维位姿和路标，以及 17728 个测量。图 3.1(a)–(c) 展示这三个问题的真值轨迹。图 3.1(d)–(f) 展示存在离群值时由最小二乘形式产生的估计。具体地，对 M3500 和 Victoria Park，我们向（回环或路标）测量中加入 15\% 的随机离群值；SubT 数据集本身已包含离群值。图中以红色显示离群测量。可以看到，离群值会导致完全错误的轨迹和地图估计。此外，离群值常揭示环境中的感知混叠：例如，SubT 数据集中间两条外观相似的竖直走廊会诱发许多伪回环，误导后端构建出只有一条竖直走廊的地图。

\begin{figure}[htbp]
\centering
\includegraphics[width=\linewidth]{01acf9ddaa14e194a185fa6e53e754763021e5e62566968febee2dbd36f5084c.jpg}
\caption{图 3.2 视觉 SLAM 问题中跨三个帧（分别采集于时刻 $k - 2$、$k - 1$ 和 $k$）的特征跟踪。内点以绿色显示：它们是随时间描绘同一（静态）三维点的像素。离群值以红色显示。}
\label{fig:01acf9ddaa14e194a185fa6e53e754763021e5e62566968febee2dbd36f5084c}
\end{figure}

\section{在 SLAM 前端检测和剔除离群值}

SLAM 前端的主要作用，是从原始传感器数据中提取中间表示或（伪）测量，之后这些量将转换为后端的因子。典型的 SLAM 前端先计算一个初始测量集合（其中可能混入许多离群值），再对该初始集合进行后处理以去除离群值。本节讨论两种在 SLAM 前端剔除离群值的方法：RANSAC 和成对一致性最大化（Pairwise Consistency Maximization）。

\section{随机采样一致性（RANSAC）}

RANSAC 是一种成熟的离群值剔除工具 \textsuperscript{[328]}，也是许多基于路标的 SLAM 系统的关键组成部分。为理解 RANSAC 是什么及其在 SLAM 中的作用，考虑一种基于路标的视觉 SLAM 方法。

\textbf{例 3.1（基于路标的视觉 SLAM 中的离群值）} 一种基于路标（或基于特征）的视觉 SLAM 方法会在每幅图像中提取二维特征点，然后通过基于光流的特征跟踪或基于描述子的特征匹配，在连续帧之间关联这些特征点（图 3.2）。具体而言，在时刻 $k$，该方法检测二维特征点，并将其与前一帧（例如时刻 $k - 1$）观测到的对应点匹配；这些匹配像素通常称为二维–二维对应（2D-2D correspondences）。由于光流或基于描述子的匹配并不精确，这一初始对应集合可能包含离群值。因此，在将其传递给估计机器人位姿和路标位置的后端之前，过滤严重离群值十分重要。

RANSAC 是一种在将对应关系传递给后端之前快速检测并去除其中离群值的工具。检测离群值依赖两个关键洞见。第一个洞见是，在 SLAM 问题中，内点对应必须满足几何约束。例如，在本例中，内点对应描述了相机运动时静态三维点的观测像素运动。由此产生的像素运动不可能任意，而必须遵循精确的几何约束，即\textbf{对极约束}（epipolar constraint）；该约束规定了两个帧中对应像素如何随相机运动发生关联。具体地，对于已标定相机，对极约束要求对应像素 $z _ { i } ( k - 1 ) , z _ { i } ( k )$，即分别在时刻 $k - 1$ 和 $k$ 描绘路标 $i$ 的像素，满足

\begin{equation}
\boldsymbol {z} _ {i} (k - 1) ^ {\top} \left([ \boldsymbol {t} _ {k} ^ {k - 1} ] _ {\times} \boldsymbol {R} _ {k} ^ {k - 1}\right) \boldsymbol {z} _ {i} (k) = 0,\tag{3.2}
\end{equation}

其中，$t _ { k } ^ { k - 1 }$ 和 $\pmb { R } _ { k } ^ { k - 1 }$ 是描述相机在时刻 $k - 1$ 与 $k$ 之间（未知）运动的相对平移和旋转。\textsuperscript{2} 更一般地，若将第 $i$ 个对应记为 $z _ { i }$（在上述例子中，$z _ { i } = \{ z _ { i } ( k - 1 ) , z _ { i } ( k ) \}$），则这类几何约束具有如下形式：

\begin{equation}
C (\boldsymbol {z} _ {i}, \boldsymbol {x}) \leq \gamma ,\tag{3.3}
\end{equation}

它表明对应关系必须满足某个不等式，而该不等式可能是未知状态 $\boldsymbol{x}$ 的函数。式（3.3）右侧的参数 $\gamma$ 通常经过调节以计及噪声的存在。例如，虽然理想情况下式（3.2）中的对极约束应严格满足，但实践中像素检测带有噪声，因而可能存在小误差；因此，我们放宽该约束，只要求 $\begin{array} { r } { | z _ { i } ( k - 1 ) ^ { \mathsf { T } } \left( [ t _ { k } ^ { k - 1 } ] _ { \times } R _ { k } ^ { k - 1 } \right) z _ { i } ( k ) \vert \leq \gamma } \end{array}$，其中 $\gamma$ 为某个较小的值。

第二个洞见是，假设离群值数量不太多，我们可以将内点找作：对某个 $x$，满足几何约束（3.3）的最大对应集合：

\begin{equation}
\begin{array}{c} \mathsf {S} _ {\mathrm{CM}} ^ {*} = \underset {\boldsymbol {x}, \mathsf {S} \subset \mathsf {M}} {\operatorname{argmax}}   | \mathsf {S} | \\ \text { s.t. } C (\boldsymbol {z} _ {i}, \boldsymbol {x}) \leq \gamma , \quad \forall i \in \mathsf {S} \end{array}\tag{3.4}
\end{equation}

其中，$M$ 是初始假定对应（putative correspondences）的集合，$|S|$ 表示子集 $S$ 的基数（元素个数）\textsuperscript{[728]}。换言之，优化问题（3.4）在假定对应集合 $M$ 中寻找最大的子集 $S$，使得 $S$ 中的测量对同一个 $x$ 都满足几何约束。直观地说，问题（3.4）刻画了如下直觉：内点（由集合 $S$ 估计）必须就同一个 $x$“达成一致”，例如它们必须全都与机器人的真实运动一致。问题（3.4）在计算机视觉中称为\textbf{一致性最大化}（consensus maximization）。注意，式（3.4）不要求求解完整的 SLAM 问题（后者可能包含许多位姿和路标），因为它只涉及 SLAM 状态的一小部分；例如，对极约束（3.2）只涉及两个帧之间的相对位姿，而非整个 SLAM 轨迹。同时，问题（3.4）仍然是一个困难的组合问题，这与典型 SLAM 前端对快速运行时间的要求相冲突。因此，人们通常不寻求式（3.4）的精确解，而是采用快速启发式方法近似地求解它。

随机采样一致性（RANdom SAmple Consensus，RANSAC）可能是为一致性最大化问题（3.4）寻找近似解的最著名方法。RANSAC 建立在一个关键假设之上：式（3.4）中的 $x$ 维数相对较低，且可利用快速估计器（即所谓的最小求解器，minimal solvers）从一小组测量（即所谓的最小集，minimal set）中估计得到。\textsuperscript{3} 例如，在我们的视觉 SLAM 示例中，可利用 Nister 的五点法 \textsuperscript{[810]}，仅通过 5 个像素对应估计两个相机帧之间的相对运动。RANSAC 的核心思想是：无需穷尽检查每一个可能的子集 $\mathsf { S } \subset \mathsf { M }$，而可以对最小测量集进行采样以寻找内点。更具体地，RANSAC 迭代执行以下三步：

1. 采样 $n$ 个对应关系的子集，其中 $n$ 是当前问题最小集的大小；\textsuperscript{4}

2. 使用最小求解器，根据采样的 $n$ 个对应关系计算估计 $\hat{\pmb{x}}$；\textsuperscript{5}

3. 对于上一步计算得到的 $\hat{\pmb{x}}$，选择满足几何约束 $C ( z _ { i } , \hat { \pmb x } ) \leq \gamma$ 的对应关系 $ \mathsf { S } \subset \mathsf { M }$。若该集合 $S$ 大于先前迭代中计算的集合，则将其保存。

最后一步计算得到的集合 $S$ 称为\textbf{一致集}（consensus set）。RANSAC 通常在计算到足够大的一致集（由用户参数指定）后，或达到最大迭代次数后停止。RANSAC 本质上试图从测量集合中采样 $n$ 个内点，因为这些内点很可能会与其他所有内点“达成一致”，因而具有较大的一致集。

RANSAC 是众多离群值剔除问题的首选方案。特别是，在离群值数量较少且最小集较小时，它能迅速收敛到良好估计（即良好的对应集合）。若从测量集合中采样一个内点的概率为 $\omega$，\textsuperscript{6} 则可容易得出，RANSAC 找到一组内点所需的期望迭代次数为 $\scriptstyle { \frac { 1 } { \omega ^ { n } } }$。例如，当 $n = 5$ 且 $\omega = 0 . 7$（即 70\% 的测量为内点）时，期望迭代次数小于 10。加之非最小求解器在实践中极快（常能在数毫秒内完成数百次迭代），RANSAC 极具吸引力。此外，RANSAC 还会提供估计 $\hat{\pmb{x}}$（例如机器人里程计），可作为后端的初始猜测。

另一方面，RANSAC 未必适合所有问题。尤其是，当内点数量很少或最小集很大时，期望迭代次数会变得不切实际地大。例如，当 $n = 1 0$ 且 $\omega = 0 . 1$ 时，找到一组内点所需的期望迭代次数达到 $1 0 ^ { 1 0 }$；在更少的迭代次数后终止 RANSAC，很可能返回错误解（即错误的 $\hat{\pmb{x}}$ 和对应关系）。如下一节所述，在许多 SLAM 问题中，“内点很多且最小集很小”的假设并不总是成立。

\section{基于图论的离群值剔除与成对一致性最大化}

如前所述，当离群值比例适中（例如低于 70\%）且最小集较小（例如小于 8）时，RANSAC 非常有效。然而，存在严重感知混叠的环境可能具有极高比例的离群值。此外，我们感兴趣的并非所有问题都存在最小集很小的快速最小求解器。例如，考虑一个具有 $N$ 个节点的位姿图 SLAM 问题，最小集至少必须包括 $N - 1$ 个测量（构成位姿图的一棵生成树），而 $N$ 通常有数千之多。

\begin{figure}[htbp]
\centering
\includegraphics[width=\linewidth]{84029a2147cbfeb042fa3a747daf672039fa3529a7e21e2ef318d1a92d5c4800.jpg}
\caption{图 3.3 (a) 来自两幅 RGB-D 扫描的三维–三维对应，它们表示场景的两个局部视图。绿色线条表示内点对应，红色线条表示离群值对应。(b) 回环中含离群值的位姿图。绿色虚线表示内点回环，红色点线表示一个离群回环。}
\label{fig:84029a2147cbfeb042fa3a747daf672039fa3529a7e21e2ef318d1a92d5c4800}
\end{figure}

因此，本节介绍一种称为\textbf{成对一致性最大化}（Pairwise Consistency Maximization，PCM）的替代方法。该方法不对最小集采样，而是借助图论寻找内部彼此“一致”的最大测量集合。该方法可处理离群值比例超过 90\% 的测量集合，并在传递给后端之前剪除严重离群值。该方法最初由文献 \textsuperscript{[728]} 提出，并在 \textsuperscript{[1005, 332]} 中推广至超越成对一致性的情况。

PCM 的关键洞见是：对于许多问题，可以定义一致性函数来描述一对测量是否相互一致。下面通过两个例子说明这一点。

\textbf{例 3.2（带 RGB-D 相机的基于路标视觉 SLAM 中的一致性函数）} 使用 RGB-D 相机的基于路标视觉 SLAM 方法会在每个 RGB-D 帧中提取三维特征点，再在连续帧之间将它们关联起来（图 3.3(a)）。具体地，在时刻 $k$，该方法检测三维特征点，并将其与前一帧（例如时刻 $k - 1$）观测到的对应点进行匹配；这些匹配的三维点通常称为三维–三维对应（3D-3D correspondences）。时刻 $k$ 和 $k - 1$ 采集的三维点，理想情况下对应于从两个不同视点观测到的同一组三维静态点；因此，对应点对 $\{ z _ { i } ( k - 1 ) , z _ { j } ( k - 1 ) \}$ 与 $\{ z _ { i } ( k ) , z _ { j } ( k ) \}$ 之间的距离必须随时间保持不变（至多存在噪声）：

\begin{equation}
\left| \left\| \boldsymbol {z} _ {i} (k - 1) - \boldsymbol {z} _ {j} (k - 1) \right\| - \left\| \boldsymbol {z} _ {i} (k) - \boldsymbol {z} _ {j} (k) \right\| \right| \leq \gamma\tag{3.5}
\end{equation}

可以看到，与 RANSAC 中使用的几何约束不同，一致性函数（3.5）(i) 不依赖于状态，因而无需最小求解器即可直接求值；(ii) 无论最小集大小如何，均涉及一对对应关系。尽管前一个例子也可使用 $\mathrm { R A N S A C }$ 求解，\textsuperscript{7} 现在考虑一个维度更高的问题。

\textbf{例 3.3（位姿图 SLAM 中的一致性函数）} 考虑一个位姿图 SLAM 问题，其中回环可能因地点识别失败和感知混叠而包含离群值；我们假设里程计可靠且不含离群值。为了判断两个回环是否彼此一致，需要注意：在无噪声情形下，图中沿环的位姿测量复合后必须为单位元（图 3.3(b)）。\textsuperscript{8} 因此，一对回环 $\pmb { T _ { b } ^ { a } }$（位姿 $a$ 与 $b$ 之间）和 $\pmb { T } _ { d } ^ { c }$（位姿 $c$ 与 $d$ 之间）必须满足：

\begin{equation}
\mathrm{dist} (\pmb {T} _ {b} ^ {a} \cdot \bar {\pmb {T}} _ {c} ^ {b} \cdot \pmb {T} _ {d} ^ {c} \cdot \bar {\pmb {T}} _ {a} ^ {d}, \mathbf {I}) \leq \gamma\tag{3.6}
\end{equation}

其中，$\bar { \pmb { T } } _ { c } ^ { b }$ 和 $\bar { \pmb { T } } _ { a } ^ { d }$ 分别是从节点 $b$ 到节点 $c$、以及从节点 $d$ 到节点 $a$ 的里程计测量链；$\mathrm{dist}$ 是合适的距离函数，用于度量 $\pmb { T } _ { b } ^ { a } \cdot \bar { \pmb { T } } _ { c } ^ { b } \cdot \pmb { T } _ { d } ^ { c} \cdot \bar { \pmb { T } } _ { a } ^ { d }$ 与单位位姿之间的距离。与往常一样，$\gamma$ 是为计及噪声而选择的参数：由于里程计和回环中存在噪声，沿一个环的测量未必复合为单位元。\textsuperscript{9}

更一般地，一致性函数是关联两个测量并且必须满足给定约束的函数。对于一对测量 $z _ { i }$ 和 $z _ { j }$，所得的成对一致性约束形式为：

\begin{equation}
F (\boldsymbol {z} _ {i}, \boldsymbol {z} _ {j}) \leq \gamma ,\tag{3.7}
\end{equation}

其中，$F$ 是一致性函数，$\gamma$ 是计及测量噪声的用户指定参数。需要强调，成对一致性约束与状态无关，因此只需检查每一对测量，无需求助于最小求解器即可高效验证。

利用式（3.7），我们可以给出另一种离群值剔除方法：选择彼此成对一致的最大测量集合：

\begin{equation}
\begin{array}{c} \mathsf {S} _ {\mathrm{PCM}} ^ {*} = \underset {\mathsf {S} \subset \mathsf {M}} {\operatorname{argmax}} | \mathsf {S} | \\ \text {s.t.} F (\boldsymbol {z} _ {i}, \boldsymbol {z} _ {j}) \leq \gamma , \quad \forall i, j \in \mathsf {S} \end{array}\tag{3.8}
\end{equation}

问题（3.8）寻找最大的测量子集 $S$，使得 $S$ 中每对测量均成对一致。我们将其称为成对一致性最大化（PCM）问题。该问题仍具有组合性质，但看起来比式（3.4）稍易：它不涉及 $x$，且可对 $M$ 中每一对 $( i , j )$ 预先计算约束 $F ( z _ { i } , z _ { j } ) \leq \gamma$。此外，该问题具有图论解释，因此可使用图论中的成熟工具，即最大团算法，来求解。

为建立问题（3.8）与图论的联系，将离群值剔除问题可视化为一个图 $\mathcal { G }$：图的节点是假定测量 $i \in { \mathsf { M } }$；若

\begin{figure}[htbp]
\centering
\includegraphics[width=\linewidth]{f405223ca19647a424ddeb6d4d8375573b5020f5b580ad455f92e0a6a438e986.jpg}
\caption{图 3.4 (a) 图 3.3(a) 中三维–三维对应示例的一致性图。(b) 图 3.3(b) 中位姿图示例的回环一致性图。}
\label{fig:f405223ca19647a424ddeb6d4d8375573b5020f5b580ad455f92e0a6a438e986}
\end{figure}

$F ( z _ { i } , z _ { j } ) \leq \gamma$，则节点 $i$ 和 $j$ 之间存在一条边。这通常称为\textbf{一致性图}（consistency graph）。问题（3.8）现在要求选择最大的节点子集 $S$，使 $S$ 中每对节点之间都由一条边相连；这正是图中\textbf{最大团}（maximum clique）的定义。更具体地，在图论中，团是一个节点子集，其中每一对节点之间都有边；最大团则是图中这样的最大节点子集。因此，问题（3.8）的解就是一致性图的最大团。这一图论联系在实践中很有价值，因为对给定图寻找最大团是图论中研究充分的问题，称为最大团问题。最大团问题是 NP-hard 问题 \textsuperscript{[1197]}，且难以近似 \textsuperscript{[1314, 322]}，这意味着寻找任意接近真实解的解同样是 NP-hard 的。不过，已有数十种最大团算法被提出，其中一些可根据图的稠密度处理规模相当大的问题。所提出的方法大多可分为精确方法与基于启发式的方法。所有精确算法的复杂度均为指数级，且通常基于分支定界；启发式算法则常试图利用问题中的某种结构而更快，但代价是不一定保证最优解 \textsuperscript{[1197]}。相对较新的工作，如 \textsuperscript{[853]}，提出了可并行化的最大团算法，能够在大型稀疏图中快速寻找最大团。

总而言之，使用最大团算法求解 PCM 问题包括以下步骤：

1. 为当前问题选择一致性函数 $F$。

2. 对每一对假定测量 $( i , j ) \in M$ 求值一致性函数；当 $F ( z _ { i } , z _ { j } ) \leq \gamma$ 时，在 $i$ 与 $j$ 之间连边，从而构建一致性图。

3. 使用精确或近似最大团算法求解一致性图的最大团。

4. 返回（可能为近似的）最大团中的测量 $S$。

需要指出，一致性函数的选择依赖于具体问题；选择好的函数会显著影响离群值剔除质量。例如，可以选择一个无意义函数 $F ( z _ { i } , z _ { j } ) = 0$，它无论输入为何都返回零；此类函数无法剔除任何离群值，从而使 PCM 无效。反之，若构造的函数使得只有内点能通过检验，则可恰好剔除所有离群值。文献 \textsuperscript{[1005, 332]} 讨论了适用于广泛几何问题的一组潜在一致性函数。

在结束本节前，还需说明几点。尽管 PCM 相较 RANSAC 能处理更高比例的离群值，并且更适合某些问题（例如位姿图 SLAM），但 PCM 与 RANSAC 的权衡更为细致。RANSAC 使用最小求解器计算的估计来评估单个测量的一致性；PCM 则以成对方式评估一组测量彼此之间的一致性。在某些情形中，RANSAC 的单测量一致性不足以评估整个测量集合；位姿图优化常是如此，因为一对回环测量各自的一致性并不必然保证这两个回环之间成对一致。\textsuperscript{10} 另一方面，对于三维–三维位姿估计（例 3.2）等某些问题，PCM 使用的成对一致性函数（3.5）可能比 RANSAC 更宽松，从而将某些离群值归类为内点。

在 PCM 的语境下，还应注意，精确最大团求解器在稠密一致性图中（即许多测量对彼此一致时）往往较慢，因此对某些问题而言，基于启发式的最大团解可能是更好的选择。最后，对某些问题很难设计合适的一致性函数；例如，对于二维–二维对应，由于缺乏合适的不变量，难以严格地设计一个通用一致性函数（参见 \textsuperscript{[1005]} 的讨论）。

\section{提高 SLAM 后端对离群值的鲁棒性}

包括 RANSAC 和 PCM 在内的前端离群值剔除仍可能漏检离群值，并向后端传递受离群值污染的测量集合。\textsuperscript{11} 如第 3.1.2 节所见，采用标准最小二乘估计时，少数离群值即可导致完全错误的结果。因此，提高后端对剩余离群值的鲁棒性十分重要。

第 3.1.2 节指出，使用残差平方会“放大”离群测量对代价函数的影响。本节遵循鲁棒统计中 M-估计（M-Estimation）的标准理论 \textsuperscript{[495]}，对 SLAM 优化目标函数作轻微修改，以重新获得对离群值的鲁棒性。

M-估计（“maximum-likelihood-type estimation”）是鲁棒估计的一个框架。它建议将式（3.1）中的平方损失替换为恰当选择的鲁棒损失函数 $\rho$：

\begin{equation}
\boldsymbol {x} ^ {\text { MAP }} = \underset {\boldsymbol {x}} {\arg \min} \sum_ {i} \boldsymbol {r} _ {i} (\boldsymbol {x}) ^ {2} \quad \Longrightarrow \quad \boldsymbol {x} ^ {\text { MEST }} = \underset {\boldsymbol {x}} {\arg \min} \sum_ {i} \rho \left(\boldsymbol {r} _ {i} (\boldsymbol {x})\right).\tag{3.9}
\end{equation}

鲁棒损失 $\rho$ 的关键要求是：它为非负函数，且对于大残差的增长速度低于二次；换言之，随着 $r _ { i }$ 变大，鲁棒损失函数需要具有导数 $\begin{array} { r } { \frac { \partial \rho ( \pmb { r } _ { i } ) } { \partial \pmb { r } _ { i } } \ll \frac { \partial \| \pmb { r } _ { i } \| ^ { 2 } } { \partial \pmb { r } _ { i } } = 2 \pmb { r } _ { i } } \end{array}$；在许多情况下，希望当 $r _ { i }$ 变大时，$\frac { \partial \rho ( \pmb { r } _ { i } ) } { \partial \pmb { r } _ { i } }$ 趋于零。为阐明这一要求，考虑使用梯度下降求解 $\min _ { \mathbf { \mathcal { x } } } \sum _ { i } \rho \left( \pmb { r } _ { i } ( \pmb { x } ) \right)$。由链式法则，目标 $f ( \pmb { x } ) \doteq \sum _ { i } \rho \left( \pmb { r } _ { i } ( \pmb { x } ) \right)$ 的梯度为：

\begin{equation}
\frac {\partial f}{\partial \boldsymbol {x}} = \sum_ {i} \frac {\partial \rho (\boldsymbol {r} _ {i})}{\partial \boldsymbol {r} _ {i}} \cdot \frac {\partial \boldsymbol {r} _ {i} (\boldsymbol {x})}{\partial \boldsymbol {x}}\tag{3.10}
\end{equation}

由式（3.10）可知，若从一个良好的初始猜测开始（即相对接近真值），离群测量具有较大残差，因而 $\frac { \partial \rho ( \pmb { r } _ { i } ) } { \partial \pmb { r } _ { i } }$ 很小，对总体下降方向的影响也很小。因此它们在估计中几乎没有影响。实际上，函数 $\begin{array} { r } { \psi ( \pmb { r } _ { i } ) : = \frac { \partial \rho ( \pmb { r } _ { i } ) } { \partial \pmb { r } _ { i } } } \end{array}$ 通常称为\textbf{影响函数}（influence function）\textsuperscript{[85]}。

鲁棒估计文献并未给出唯一的鲁棒损失函数，而是提供了多种可能选择。图 3.5 列出了常见损失函数，其中包括 Huber、Geman-McClure、Tukey 双权（Tukey's biweight）和截断二次损失等常用鲁棒损失，也包括一种更激进的选择，即最大一致性损失。后一种通常不列在鲁棒估计文献的损失函数之中，但这里提及它，是因为它与我们在式（3.4）中讨论的一致性最大化问题相联系。\textsuperscript{12} 鲁棒损失的选择高度依赖具体问题 \textsuperscript{[554, 1075]}。例如，当已知合理的截断阈值（即内点的预期最大误差）时，具有硬截断的损失函数，例如截断二次损失（其二次区与函数“平坦”部分之间发生突变），更为合适。还必须考虑计算因素。例如，Huber 常用于 BA 问题，因为它是凸函数，在优化过程中表现更稳定，尽管它对离群值保留非零影响（只有当大残差时损失保持常数，影响才变为零）。另一方面，截断二次和最大一致性损失对离群值尤不敏感，但常需要专用求解器。\textsuperscript{13}

\begin{figure}[htbp]
\centering
\includegraphics[width=\linewidth]{d62f0b14a781770de0d928a2bc5db48feb73bad0b5cecf4fe695e8f16d77a901.jpg}
\label{fig:d62f0b14a781770de0d928a2bc5db48feb73bad0b5cecf4fe695e8f16d77a901}
\end{figure}

(a) 二次损失

\begin{figure}[htbp]
\centering
\includegraphics[width=\linewidth]{c57bf228a0cde1e02ee2cc1c508ef0c788bff27c89ec49502de7fb1caba5964f.jpg}
\label{fig:c57bf228a0cde1e02ee2cc1c508ef0c788bff27c89ec49502de7fb1caba5964f}
\end{figure}

(b) Huber 损失

\begin{figure}[htbp]
\centering
\includegraphics[width=\linewidth]{abd3380a85d6de0464cb838eddd7048f4f923ca178f7f16256765e143810ad3b.jpg}
\label{fig:abd3380a85d6de0464cb838eddd7048f4f923ca178f7f16256765e143810ad3b}
\end{figure}

(c) Geman-McClure

\begin{figure}[htbp]
\centering
\includegraphics[width=\linewidth]{f41954a5877f37fd3f5089ab6964361b69b75d4b3ec9ee80360e1506a2277964.jpg}
\label{fig:f41954a5877f37fd3f5089ab6964361b69b75d4b3ec9ee80360e1506a2277964}
\end{figure}

(d) Tukey 双权

\begin{figure}[htbp]
\centering
\includegraphics[width=\linewidth]{f458be9aab9f793c517fb2748ea55d4d143c45f62830ce2e66af3492680a7929.jpg}
\label{fig:f458be9aab9f793c517fb2748ea55d4d143c45f62830ce2e66af3492680a7929}
\end{figure}

\begin{figure}[htbp]
\centering
\includegraphics[width=\linewidth]{0ba2757187eef46948d72a72d4db2b40f92e9dab77d541d37f5f9517c1f5d87c.jpg}
\label{fig:0ba2757187eef46948d72a72d4db2b40f92e9dab77d541d37f5f9517c1f5d87c}
\end{figure}

(e) 截断二次
(f) 最大一致性损失

图 3.5 二次损失与鲁棒损失函数示例。鲁棒损失函数的形状由一个控制内点和离群值分界的参数确定。

图 3.6(g)–(l) 展示了对上述两种鲁棒损失，即 Huber 损失和截断二次损失，应用梯度下降所得到的 SLAM 轨迹。这里采用与图 3.1 相同的数据集。我们使用启用 gradientDescent 标志的 GTSAM NonlinearConjugateGradientOptimizer \textsuperscript{[253]} 实现梯度下降求解器，并使用鲁棒噪声模型实例化鲁棒损失函数。最大迭代次数和停止条件阈值（相对与绝对容差）分别设为 10000 和 $1 0 ^ { - 7 }$，其余参数保留为 GTSAM 默认值。图中，若一条边被优化正确归类为内点或离群值，则显示为灰色（即落入 Huber 或截断二次损失二次区域的内点）；若离群值被错误归为内点，即“假阳性”（false positive），则显示为红色；若内点被错误归为离群值，即“假阴性”（false negative），则显示为蓝色。与非鲁棒最小二乘优化（图 3.1(d)–(f)）相比，鲁棒损失可使 M3500 和 SubT 数据集迅速恢复离群值鲁棒性；但简单的梯度下降仍可能无法正确优化截断二次代价等高度非凸函数，如 Victoria Park 数据集所示。下文将用基于渐进非凸性（graduated non-convexity）的更好求解器处理该问题。此外，尽管梯度下降在许多实例中已改善性能（图 3.6），它仍具有缓慢的收敛尾部；在我们的实验中，往往需要数千次迭代才收敛。因此，本节其余部分讨论同时提高收敛质量和速度的更高级求解器。

M3500
\begin{figure}[htbp]
\centering
\includegraphics[width=\linewidth]{94b8f7648085474b55ce2dca350bb3e76fd2834866375f7f01852ee6305aaca1.jpg}
\label{fig:94b8f7648085474b55ce2dca350bb3e76fd2834866375f7f01852ee6305aaca1}
\end{figure}

(a)
SubT

\begin{figure}[htbp]
\centering
\includegraphics[width=\linewidth]{895b8b8f94a3ce2bb52899e3677c80b447412234411ba1eb8645019aa617a9cb.jpg}
\label{fig:895b8b8f94a3ce2bb52899e3677c80b447412234411ba1eb8645019aa617a9cb}
\end{figure}

(b)
Victoria Park

\begin{figure}[htbp]
\centering
\includegraphics[width=\linewidth]{0fe34b30ed9dd12ac0a191a3c98a1553402d66ff07694f212c4847e7fd0ffb11.jpg}
\label{fig:0fe34b30ed9dd12ac0a191a3c98a1553402d66ff07694f212c4847e7fd0ffb11}
\end{figure}

\begin{figure}[htbp]
\centering
\includegraphics[width=\linewidth]{740e4b3e024d267da11078510075d736709ab6c40b47a8a281810f6c36e90088.jpg}
\label{fig:740e4b3e024d267da11078510075d736709ab6c40b47a8a281810f6c36e90088}
\end{figure}

(g)

(c)
\begin{figure}[htbp]
\centering
\includegraphics[width=\linewidth]{e495738bc6f276a9f347519cec47034b22a623e4cf349d459d1ad90c61e1d976.jpg}
\label{fig:e495738bc6f276a9f347519cec47034b22a623e4cf349d459d1ad90c61e1d976}
\end{figure}

\begin{figure}[htbp]
\centering
\includegraphics[width=\linewidth]{d60fedf39af3e05f868b32152c41a5f32aac657c9804c48482bdb46bb5c314bd.jpg}
\label{fig:d60fedf39af3e05f868b32152c41a5f32aac657c9804c48482bdb46bb5c314bd}
\end{figure}

\begin{figure}[htbp]
\centering
\includegraphics[width=\linewidth]{85c8e422b6bfba361c12de6e5d4625705cbad67c27e1288e7d585a3831ea9605.jpg}
\label{fig:85c8e422b6bfba361c12de6e5d4625705cbad67c27e1288e7d585a3831ea9605}
\end{figure}

(j)

(h)
(i)
\begin{figure}[htbp]
\centering
\includegraphics[width=\linewidth]{3d959690ff503efd6feab0e359665c47153cc67e22401e00bf11f8da991c418c.jpg}
\label{fig:3d959690ff503efd6feab0e359665c47153cc67e22401e00bf11f8da991c418c}
\end{figure}

(k)

\begin{figure}[htbp]
\centering
\includegraphics[width=\linewidth]{315a9c1c1ab8526b9fb06c328de349f4eeccd799161aa0d486bc8260e16faf35.jpg}
\label{fig:315a9c1c1ab8526b9fb06c328de349f4eeccd799161aa0d486bc8260e16faf35}
\end{figure}

(l)
图 3.6 使用鲁棒损失函数和梯度下降求解含离群值 SLAM 问题：(a)–(c) 为 M3500、SubT 和 Victoria Park 数据集的真值轨迹；(d)–(f) 为使用梯度下降和 Huber 损失得到的轨迹估计；(g)–(l) 为使用梯度下降和截断二次损失得到的轨迹估计。测量以彩色边显示：灰色边为被优化正确归类为内点或离群值的边；红色边为被错误归类为内点的离群值（假阳性）；蓝色边为被错误归类为离群值的内点（假阴性）。

在深入讨论高级求解器前，最后指出：虽然切换到鲁棒损失函数看似放弃了概率框架，实际上可通过对重尾噪声分布应用 MAP 估计推导出若干鲁棒损失。例如，当假设噪声服从描述内点的高斯密度与描述离群值的均匀分布之间的 max-mixture 分布时，MAP 估计会导出截断二次损失 \textsuperscript{[34]}。

\section{迭代重加权最小二乘}

M-估计在式（3.9）中用鲁棒损失 $\rho$ 替换最小二乘损失，即用一个对大残差次二次增长的函数替换它。这有两个代价。首先，我们失去了为最小二乘形式开发的高效解法；例如，高斯–牛顿和 Levenberg–Marquardt 方法是为最小二乘问题设计的。其次，由于 M-估计通常具有非凸的代价景观，迭代求解器（例如基于梯度下降的方法）对初始化质量敏感，常收敛至不期望的次优估计。本节介绍一种求解 M-估计的常用算法，即\textbf{迭代重加权最小二乘}（iteratively reweighted least squares，IRLS）；顾名思义，它允许复用高效的最小二乘求解器。下一节将介绍渐进非凸性，以改善 IRLS 的收敛性。

IRLS 的基本思想是：在每次迭代中通过求解加权最小二乘问题来优化式（3.9）：

\begin{equation}
\boldsymbol {x} ^ {(t + 1)} = \underset {\boldsymbol {x}} {\arg \min} \sum_ {i} w _ {i} (\boldsymbol {x} ^ {(t)}) r _ {i} ^ {2} (\boldsymbol {x}),\tag{3.11}
\end{equation}

其中，权重 $w _ { i }$ 取决于上一次迭代的估计 $\mathbf { \boldsymbol { x } } ^ { ( t ) }$。我们希望迭代解 $\mathbf { \boldsymbol { x } } ^ { ( t ) }$ 收敛到 M-估计（3.9）的最优解。这意味着，式（3.10）所示鲁棒损失 $\rho$ 的梯度，必须与式（3.11）中损失的梯度相匹配。将式（3.11）的梯度写为

\begin{equation*}
\sum_ {i} 2 w _ {i} (\boldsymbol {x} ^ {(t)}) r _ {i} (\boldsymbol {x}) \frac {\partial r _ {i} (\boldsymbol {x})}{\partial \boldsymbol {x}}
\end{equation*}

并与式（3.10）比较，可得 IRLS 权重更新规则：

\begin{equation}
w _ {i} (\pmb {x} ^ {(t)}) = \frac {1}{2 r _ {i} (\pmb {x} ^ {(t)})} \frac {\partial \rho (r _ {i} (\pmb {x} ^ {(t)}))}{\partial r _ {i} (\pmb {x} ^ {(t)})} = \frac {\psi (r _ {i} (\pmb {x} ^ {(t)}))}{2 r _ {i} (\pmb {x} ^ {(t)})},\tag{3.12}
\end{equation}

其中，$\begin{array} { r } { \psi ( r _ { i } ) : = \frac { \partial \rho ( r _ { i } ) } { \partial r _ { i } } } \end{array}$ 是影响函数。因而，IRLS 在为每个测量 $i$ 计算权重 $w _ { i } ( \pmb { x } ^ { ( t ) } )$ 与对加权最小二乘问题（3.11）执行优化步骤（即一次 Gauss–Newton 或 Levenberg–Marquardt 迭代）之间交替进行。

图 3.7 给出了 IRLS 在 M3500、SubT 和 Victoria Park 数据集上的表现。IRLS 在数十次迭代内收敛，通常比梯度下降快得多；例如，在我们的 M3500 实验中，梯度下降优化 Huber 损失约需 5 秒，而 IRLS 少于 1.5 秒。另一方面，这种更快的收敛往往以精度略有下降为代价，比较图 3.7 与图 3.6 即可看出。文献 \textsuperscript{[10, 817]} 已研究了更新规则（3.12）的收敛性质。

\section{Black–Rangarajan 对偶性}

权重更新规则（3.12）在实践中广泛使用，但其推导带有一定启发性；此外，在 $\rho$ 的不可微点（例如截断二次损失的截断点）处，式（3.12）没有定义。下面引入一个更有原则的框架，即 Black–Rangarajan（B–R）对偶性 \textsuperscript{[85]}，用于通过 IRLS 求解 M-估计。

使用截断二次损失说明 B–R 对偶性的直觉：

\begin{equation}
\rho (r _ {i} (\boldsymbol {x})) := \min \{r _ {i} ^ {2} (\boldsymbol {x}), \beta_ {i} ^ {2} \},\tag{3.13}
\end{equation}

其中，$\beta _ { i } ^ { 2 }$ 是第 $i$ 个残差的界：若 $r _ { i } ^ { 2} ( { \pmb x } ) \leq \beta _ { i } ^ { 2}$，第 $i$ 个测量为内点，否则为离群值。通过引入新的权重变量 $w _ { i } \in [ 0 , 1 ]$，可将式（3.13）的代价等价地写作两项之和：

\begin{equation}
\rho (r _ {i} (\boldsymbol {x})) := \min _ {w _ {i} \in [ 0, 1 ]} w _ {i} r _ {i} ^ {2} (\boldsymbol {x}) + (1 - w _ {i}) \beta_ {i} ^ {2},\tag{3.14}
\end{equation}

其中第一项恰为加权最小二乘，第二项是与 $x$ 无关的 $w _ { i }$ 的函数。借助式（3.14），使用截断二次损失的 M-估计问题（3.9）可改写为

\begin{equation}
\min _ {\boldsymbol {x},\, w _ {i} \in [ 0, 1 ],\, i = 1, \dots , N} \sum_ {i} \left[ w _ {i} r _ {i} ^{2} (\boldsymbol {x}) + (1 - w _ {i}) \beta_ {i} ^ {2} \right],\tag{3.15}
\end{equation}

其中，对每个测量残差 $r _ { i } ( { \pmb x } )$ 都引入一个 $w _ { i }$。问题（3.15）容易解释：$w _ { i } = 1$ 意味着 $r _ { i } ^ { 2} ( { \pmb x } ) \leq \beta _ { i } ^ { 2}$，第 $i$ 个测量为内点；$w _ { i } = 0$ 意味着 $r _ { i } ^ { 2} ( { \pmb x } ) > \beta _ { i } ^ { 2}$，第 $i$ 个测量为离群值。此外，所有 $w _ { i } = 0$ 的残差均被有效地从优化问题（3.15）中丢弃，因而保证了鲁棒性。

B–R 对偶性将上述推导推广到一族鲁棒损失。

\textbf{定理 3.4（Black–Rangarajan 对偶性 \textsuperscript{[85]}）} 给定鲁棒损失函数 $\rho ( \cdot )$，定义 $\phi ( z ) : = \rho ( \sqrt { z } )$。若 $\phi(z)$ 满足 $\lim _ {z \to 0} \phi ^ { \prime } ( z ) = 1$、$\lim _ {z \to \infty} \phi ^ { \prime } ( z ) = 0$ 且 $\phi ^ { \prime \prime } ( z ) < 0$，则 M-估计问题（3.9）等价于

\begin{equation}
\min _ {\boldsymbol {x},\, w _ {i} \in [ 0, 1 ],\, i = 1, \dots , N} \sum_ {i = 1} ^ {N} \left[ w _ {i} r _ {i} ^ {2} (\boldsymbol {x}) + \Phi_ {\rho} (w _ {i}) \right],\tag{3.16}
\end{equation}

其中，$w _ { i } \in [ 0 , 1 ] , i = 1 , \ldots , N$ 是与每个残差 $r _ { i }$ 关联的权重变量；函数 $\Phi _ { \rho } ( w _ { i } )$ 称为\textbf{离群值过程}（outlier process），它定义对权重 $w _ { i }$ 的惩罚，其形式取决于鲁棒损失 $\rho$ 的选择。

\begin{figure}[htbp]
\centering
\includegraphics[width=\linewidth]{1d130e48ce5ae89eb8ea8000207184c0f46a92d0179c5bd4cd972e0e41ef7401.jpg}
\label{fig:1d130e48ce5ae89eb8ea8000207184c0f46a92d0179c5bd4cd972e0e41ef7401}
\end{figure}

(a)

\begin{figure}[htbp]
\centering
\includegraphics[width=\linewidth]{22004b752ab47b2f04d0b7a15f4b9578f0daee5a9933cddbe2078e905e5b17f8.jpg}
\label{fig:22004b752ab47b2f04d0b7a15f4b9578f0daee5a9933cddbe2078e905e5b17f8}
\end{figure}

(b)

\begin{figure}[htbp]
\centering
\includegraphics[width=\linewidth]{caa9b3945f4dbad245aac6253ee8d01d1bde3bc8927179457a11efc3014a54aa.jpg}
\label{fig:caa9b3945f4dbad245aac6253ee8d01d1bde3bc8927179457a11efc3014a54aa}
\end{figure}

(c)

\begin{figure}[htbp]
\centering
\includegraphics[width=\linewidth]{efd597cb0e6be3d3c4e7ca6490473b66fb6a407d9a225458dc8e26b0cecc5a7e.jpg}
\label{fig:efd597cb0e6be3d3c4e7ca6490473b66fb6a407d9a225458dc8e26b0cecc5a7e}
\end{figure}

(d)

\begin{figure}[htbp]
\centering
\includegraphics[width=\linewidth]{28bb557ca0cd61c56726206a0b9c994e3f1b56c337e89bb31e91beddca632adb.jpg}
\label{fig:28bb557ca0cd61c56726206a0b9c994e3f1b56c337e89bb31e91beddca632adb}
\end{figure}

(e)

\begin{figure}[htbp]
\centering
\includegraphics[width=\linewidth]{fbbe5ebf8b01008f045c9ff6c693d5f1b34e80aabd42b74778dbf831b1ea693b.jpg}
\label{fig:fbbe5ebf8b01008f045c9ff6c693d5f1b34e80aabd42b74778dbf831b1ea693b}
\end{figure}

\begin{figure}[htbp]
\centering
\includegraphics[width=\linewidth]{38194e893d58a0b04f16405f4ebcd6be594f23bae01374191eecebd6b9c446ec.jpg}
\label{fig:38194e893d58a0b04f16405f4ebcd6be594f23bae01374191eecebd6b9c446ec}
\end{figure}

\begin{figure}[htbp]
\centering
\includegraphics[width=\linewidth]{df4331d7cfcda9d6a9b97f77d0a154cedc575e3ddcc79dc0e110c75705e22a05.jpg}
\label{fig:df4331d7cfcda9d6a9b97f77d0a154cedc575e3ddcc79dc0e110c75705e22a05}
\end{figure}

(f)
\begin{figure}[htbp]
\centering
\includegraphics[width=\linewidth]{4719d820b655c27cc77f2a77b460607f77aeb990430e4ad9caa17e68cbc36e75.jpg}
\label{fig:4719d820b655c27cc77f2a77b460607f77aeb990430e4ad9caa17e68cbc36e75}
\end{figure}

(g)
\begin{figure}[htbp]
\centering
\includegraphics[width=\linewidth]{47f5bc57cbecd5a73c7edad8f51de5a1c2a3faeb922a964508678a91230f8943.jpg}
\label{fig:47f5bc57cbecd5a73c7edad8f51de5a1c2a3faeb922a964508678a91230f8943}
\end{figure}

(j)

(h)
\begin{figure}[htbp]
\centering
\includegraphics[width=\linewidth]{6adeb320e916550af9d861e27c0b5cf00b8a383c510ede33cbb133457c855d17.jpg}
\label{fig:6adeb320e916550af9d861e27c0b5cf00b8a383c510ede33cbb133457c855d17}
\end{figure}

(k)

(i)
\begin{figure}[htbp]
\centering
\includegraphics[width=\linewidth]{d395015818cefe925401bd4aecbd5f6b186796d121313b9dafefd02772df11f2.jpg}
\label{fig:d395015818cefe925401bd4aecbd5f6b186796d121313b9dafefd02772df11f2}
\end{figure}

(l)
图 3.7 使用鲁棒损失函数和迭代重加权最小二乘（IRLS）求解含离群值的 SLAM 问题：(a)–(c) 为 M3500、SubT 和 Victoria Park 数据集的真值轨迹；(d)–(f) 为使用 Huber 损失得到的轨迹估计；(g)–(i) 为使用 Geman–McClure 损失得到的轨迹估计；(j)–(l) 为使用截断二次损失得到的轨迹估计。测量以彩色边显示：灰色边为正确分类的边；红色边为被错误归类为内点的离群值（假阳性）；蓝色边为被错误归类为离群值的内点（假阴性）。

当 $\rho$ 为截断二次损失时，我们已由式（3.14）轻易得到 $\Phi _ { \rho } ( w _ { i } ) = ( 1 - w _ { i } ) \beta _ { i } ^ { 2}$。当 $\rho$ 取其他形式时，文献 \textsuperscript{[85]} 提供了推导 $\Phi _ { \rho } ( w _ { i } )$ 的方法。下面给出 Geman–McClure（G–M）鲁棒损失的例子。

\textbf{例 3.5（G–M 损失的 B–R 对偶性）} 考虑 G–M 鲁棒损失函数

\begin{equation}
\rho (r _ {i} (\pmb {x})) = \frac {\beta_ {i} ^ {2} r _ {i} ^ {2} (\pmb {x})}{\beta_ {i} ^ {2} + r _ {i} ^ {2} (\pmb {x})},\tag{3.17}
\end{equation}

其中，$\beta _ { i } ^ { 2 }$ 是与式（3.13）类似的第 $i$ 个残差的噪声界。与式（3.17）关联的离群值过程为

\begin{equation}
\Phi_ {\rho} (w _ {i}) = \beta_ {i} ^ {2} (\sqrt {w _ {i}} - 1) ^ {2}.\tag{3.18}
\end{equation}

为验证式（3.18）的正确性，考虑

\begin{equation}
\min _ {w _ {i} \in [ 0, 1 ]} w _ {i} r _ {i} ^ {2} (\boldsymbol {x}) + \Phi_ {\rho} (w _ {i}),\tag{3.19}
\end{equation}

其最优解（令式（3.19）的梯度为零）是

\begin{equation}
w _ {i} ^ {\star} = \left(\frac {\beta_ {i} ^ {2}}{r _ {i} ^ {2} (\pmb {x}) + \beta_ {i} ^ {2}}\right) ^ {2}.\tag{3.20}
\end{equation}

将式（3.20）代回式（3.19）的目标函数，即可恢复 G–M 鲁棒损失（3.17）。

\section{交替最小化}

引入 B–R 对偶性后，使用称为\textbf{交替最小化}（alternating minimization）的通用优化策略 \textsuperscript{[1112, 80]}，IRLS 算法便可自然导出。其思想是：尽管在式（3.16）中同时优化 $x$ 和 $w _ { i } \in [ 0 , 1 ]$，$i = 1 , \ldots , N$ 很困难，但固定其中一者而优化另一者很容易。确实，当 $w _ { i }$ 固定时，问题（3.16）成为加权最小二乘；类似地，当 $x$ 固定时，问题（3.16）成为

\begin{equation*}
\min _ {w _ {i} \in [ 0, 1 ], i = 1, \dots , N} \sum_ {i = 1} ^ {N} \Phi_ {\rho} (w _ {i}) + w _ {i} r _ {i} ^ {2} (\boldsymbol {x}),
\end{equation*}

它可拆为 $N$ 个子问题，每个子问题优化一个标量 $w _ { i }$：

\begin{equation}
\min _ {w _ {i} \in [ 0, 1 ]} \Phi_ {\rho} (w _ {i}) + w _ {i} r _ {i} ^ {2} (\boldsymbol {x}).\tag{3.21}
\end{equation}

问题（3.21）易于求解，且常有闭式解。事实上，对 G–M 鲁棒损失，式（3.21）的解正是式（3.20）。对截断二次损失，问题（3.21）为

\begin{equation*}
\min _ {w _ {i} \in [ 0, 1 ]} (1 - w _ {i}) \beta_ {i} ^ {2} + w _ {i} r _ {i} ^ {2} (\boldsymbol {x})
\end{equation*}

并具有闭式解

\begin{equation*}
w _ {i} ^ {\star} = \begin{cases} 1, & \text{if } r _ {i} ^ {2} (\boldsymbol {x}) < \beta _ {i} ^ {2}, \\ 0, & \text{if } r _ {i} ^ {2} (\boldsymbol {x}) > \beta _ {i} ^ {2}, \\ [0, 1], & \text{otherwise}. \end{cases}
\end{equation*}

总之，在 B–R 对偶性的语境下，第 $t$ 次 IRLS 迭代在以下两步之间交替：

1. \textbf{变量更新：} 使用当前权重 $w _ { i } ^ { ( t ) }$ 求解加权最小二乘问题

\begin{equation}
\boldsymbol {x} ^ {(t)} \in \underset {\boldsymbol {x}} {\arg \min} \sum_ {i} w _ {i} ^ {(t)} r _ {i} ^ {2} (\boldsymbol {x}).\tag{3.22}
\end{equation}

2. \textbf{权重更新：} 使用 $\mathbf { \boldsymbol { x } } ^ { ( t ) }$ 更新权重

\begin{equation}
w _ {i} ^ {(t + 1)} \in \underset {w _ {i} \in [ 0, 1 ]} {\arg \min} \Phi_ {\rho} (w _ {i}) + w _ {i} r _ {i} ^ {2} (\boldsymbol {x} ^ {(t)}), i = 1, \ldots , N.\tag{3.23}
\end{equation}

通过 B–R 对偶性得到的权重更新规则实际上与广为使用的规则（3.12）一致。有趣的是，在 SLAM 中用 G–M 鲁棒损失实例化上述 IRLS 算法，会导出动态协方差缩放（dynamic covariance scaling）算法 \textsuperscript{[11]}，该算法已在离群值鲁棒 SLAM 的语境下提出。

\section{渐进非凸性}

上一节利用 Black–Rangarajan 对偶性与交替最小化，导出了在式（3.22）和（3.23）之间交替求解的 IRLS 框架。然而，由于常见鲁棒损失的非凸性，IRLS 的收敛可能对初始化质量高度敏感，即 $\mathbf { x } ^ { ( 0 ) }$ 与式（3.9）最优解的接近程度，或 $w _ { i } ^ { ( 0 ) }$ 对各测量内点/离群值归属的刻画程度。例如，文献 \textsuperscript{[1006]} 表明，即使测量中仅有 10\% 离群值，采用截断二次损失或 Geman–McClure 损失的 IRLS 也可能失败（亦可参见图 3.7 的结果）。

本节介绍\textbf{渐进非凸性}（Graduated Non-Convexity，GNC）算法，它可使 IRLS 对初始化质量显著不那么敏感。给定鲁棒代价函数 $\rho$，GNC 的基本思想是利用标量平滑因子 $\mu$ 构造 $\rho$ 的平滑版本 $\rho _ { \mu }$。调节 $\mu$ 可控制 $\rho _ { \mu }$ 的非凸程度：在参数范围的一端，$\rho _ { \mu }$ 是凸的；在另一端，它恢复原始的 $\rho$。下面以两个例子说明。

\textbf{例 3.6（GNC 截断二次损失）} 考虑 GNC 截断二次损失函数

\begin{equation}
\rho_ {\mu} (r _ {i} (\pmb {x})) = \left\{ \begin{array}{l l} r _ {i} ^ {2} (\pmb {x}) & \mathrm{if} r _ {i} ^ {2} (\pmb {x}) \in \left[ 0, \frac {\mu}{\mu + 1} \beta_ {i} ^ {2} \right] \\ 2 \beta_ {i} | r _ {i} (\pmb {x}) | \sqrt {\mu (\mu + 1)} - \mu (\beta_ {i} ^ {2} + r _ {i} ^ {2} (\pmb {x})) & \mathrm{if} r _ {i} ^ {2} (\pmb {x}) \in \left[ \frac {\mu}{\mu + 1} \beta_ {i} ^ {2}, \frac {\mu + 1}{\mu} \beta_ {i} ^ {2} \right] \\ \beta_ {i} ^ {2} & r _ {i} ^ {2} (\pmb {x}) \in \left[ \frac {\mu + 1}{\mu} \beta_ {i} ^ {2}, + \infty \right]. \end{array} \right.\tag{3.24}
\end{equation}

当 $\mu$ 趋近于零时，$\rho _ { \mu }$ 为凸函数；当 $\mu$ 趋近于无穷大时，它恢复式（3.13）的截断二次损失。

\textbf{例 3.7（GNC Geman–McClure 损失）} 考虑 GNC Geman–McClure 损失函数

\begin{equation}
\rho_ {\mu} (r _ {i} (\pmb {x})) = \frac {\mu \beta_ {i} ^ {2} r _ {i} ^ {2} (\pmb {x})}{\mu \beta_ {i} ^ {2} + r _ {i} ^ {2} (\pmb {x})}.\tag{3.25}
\end{equation}

当 $\mu$ 趋近于无穷大时，$\rho _ { \mu }$ 为凸函数；当 $\mu = 1$ 时，它恢复 G–M 损失（3.17）。

\begin{figure}[htbp]
\centering
\includegraphics[width=\linewidth]{d5c45e32a41fb6e784e6325650985c2586a323d0397f6ad629c99100d34d40ba.jpg}
\label{fig:d5c45e32a41fb6e784e6325650985c2586a323d0397f6ad629c99100d34d40ba}
\end{figure}

(a)

\begin{figure}[htbp]
\centering
\includegraphics[width=\linewidth]{601276569c745a3004affebd09355ebb79087d71f04fbcfc5891d0779d5a58c8.jpg}
\label{fig:601276569c745a3004affebd09355ebb79087d71f04fbcfc5891d0779d5a58c8}
\end{figure}

(b)
图 3.8 使用控制参数 $\mu$ 的渐进非凸性：(a) 截断二次损失；(b) Geman–McClure 损失。\textsuperscript{[1222]}（©2020 IEEE）

图 3.8(a) 和 (b) 分别绘制了 GNC 截断二次损失和 GNC Geman–McClure 损失。可以看到，增大或减小控制参数 $\mu$ 会为函数引入更多非凸性。

式（3.24）和（3.25）中平滑 GNC 函数的一个良好性质是，B–R 对偶性仍然适用。对 GNC 截断二次损失（3.24），应用 B–R 对偶性得到离群值过程

\begin{equation*}
\Phi_ {\rho_ {\mu}} (w _ {i}) = \frac {\mu (1 - w _ {i})}{\mu + w _ {i}} \beta_ {i} ^ {2}.
\end{equation*}

对 GNC Geman–McClure 损失（3.25），应用 B–R 对偶性得到离群值过程

\begin{equation*}
\Phi_ {\rho_ {\mu}} (w _ {i}) = \mu \beta_ {i} ^ {2} (\sqrt {w _ {i}} - 1) ^ {2}.
\end{equation*}

现在可以给出 GNC 算法。每次迭代执行以下三步：

1. \textbf{变量更新：} 使用当前权重 $w _ { i } ^ { ( t ) }$ 求解加权最小二乘问题

\begin{equation}
\boldsymbol {x} ^ {(t)} \in \arg \min _ {\boldsymbol {x}} \sum_ {i} w _ {i} ^ {(t)} r _ {i} ^ {2} (\boldsymbol {x}).\tag{3.26}
\end{equation}

2. \textbf{权重更新：} 使用 $\mathbf { \boldsymbol { x } } ^ { ( t ) }$ 更新权重

\begin{equation}
w _ {i} ^ {(t + 1)} \in \underset {w _ {i} \in [ 0, 1 ]} {\arg \min} \Phi_ {\rho_ {\mu}} (w _ {i}) + w _ {i} r _ {i} ^ {2} (\boldsymbol {x} ^ {(t)}), i = 1, \ldots , N.\tag{3.27}
\end{equation}

3. \textbf{控制参数更新：} 增大或减小 $\mu$，以向 $\rho _ { \mu }$ 加入更多非凸性。

GNC 算法类似 IRLS，区别在于它从凸的平滑函数 $\rho _ { \mu }$ 开始，然后迭代更新控制参数 $\mu$，逐渐向 $\rho _ { \mu }$ 加入更多非凸性，以逼近原始损失函数 $\rho$。依据平滑损失 $\rho _ { \mu }$ 的定义，通过增大或减小 $\mu$ 可恢复原始 $\rho$。例如，GNC 截断二次损失在 $\mu$ 较大时恢复原始截断二次损失，因此每次 GNC 迭代将 $\mu$ 乘以常数因子 $\gamma > 1$（例如 \textsuperscript{[1222]} 中的 $\gamma = 1 . 4$）。反之，GNC Geman–McClure 损失在 $\mu$ 接近 1 时恢复原始 G–M 损失，因此每次 GNC 迭代将 $\mu$ 除以 $\gamma > 1$。

图 3.9 展示在图 3.6 的相同三个数据集上应用 GNC 所得的 SLAM 轨迹，采用两种不同的鲁棒损失：Geman–McClure 损失和截断二次损失。我们用 GTSAM 的 GNCOptimizer 实现 GNC。比较图 3.6 和图 3.7 可知，GNC 保证了更好的收敛性（即较不易陷入局部极小值），并在全部三个数据集中恢复出相当准确的轨迹。尽管 GNC 已表现出对离群值极强的耐受性（例如，在真实问题中可耐受约 80\%–90\% 的错误回环 \textsuperscript{[1222, 165]}），但该方法不提供收敛保证。此外，实验表明其性能依赖具体问题；它在位姿图优化问题中表现优异，但在其他感知问题中性能会显著退化 \textsuperscript{[1005]}。

下面展示一个 GNC 失败的问题，并说明组合前端和后端离群值剔除可能是有益的。为此，考虑一个比此前更具挑战性的 SLAM 设置。本章前面讨论位姿图优化问题（例如图 3.1）时，假设里程计可靠，而回环或路标测量中可能存在离群值。该设置本质上假定存在一条“里程计骨架”（odometry backbone）：它提供一组可信测量，并可构建机器人位姿的初始猜测。虽然这一假设在许多问题中现实可行，\textsuperscript{14} 但某些 SLAM 应用可能缺乏里程计骨架。例如，由于轮式里程计中的车轮打滑或 LiDAR 里程计中的错误对齐，某些里程计传感器可能产生错误测量。另一个例子是多机器人 SLAM：每个机器人都有里程计骨架，但包括全部机器人轨迹的整体 SLAM 问题却没有 \textsuperscript{[728]}。

M3500
\begin{figure}[htbp]
\centering
\includegraphics[width=\linewidth]{b05ceda71ff0e3dee01d0b5fc5336292fd937615309e129bd5b19a6083a6bc62.jpg}
\label{fig:b05ceda71ff0e3dee01d0b5fc5336292fd937615309e129bd5b19a6083a6bc62}
\end{figure}

(a)
SubT

\begin{figure}[htbp]
\centering
\includegraphics[width=\linewidth]{2509e7220eb57b3d6e615fb69bd6ab471a68341f555208c627fe8c275c27d372.jpg}
\label{fig:2509e7220eb57b3d6e615fb69bd6ab471a68341f555208c627fe8c275c27d372}
\end{figure}

Victoria Park

\begin{figure}[htbp]
\centering
\includegraphics[width=\linewidth]{cae6574d230e4a3c126702533cde19830151b1098536973603ff71a249b25185.jpg}
\label{fig:cae6574d230e4a3c126702533cde19830151b1098536973603ff71a249b25185}
\end{figure}

(b)
(c)

\begin{figure}[htbp]
\centering
\includegraphics[width=\linewidth]{5f3d80b8c9993f563c9cce32d7c19178eeaba68ef463a812f5e44b2b4eaee428.jpg}
\label{fig:5f3d80b8c9993f563c9cce32d7c19178eeaba68ef463a812f5e44b2b4eaee428}
\end{figure}

(d)

\begin{figure}[htbp]
\centering
\includegraphics[width=\linewidth]{a47812bf602bcbfc1e8bea276ffc29ff1c887a40d4a7f1088be0855292b44bf1.jpg}
\label{fig:a47812bf602bcbfc1e8bea276ffc29ff1c887a40d4a7f1088be0855292b44bf1}
\end{figure}

\begin{figure}[htbp]
\centering
\includegraphics[width=\linewidth]{def2c00b9208c3e958416444ea23922ce5d260c838da2437dadb5e5f110d6877.jpg}
\label{fig:def2c00b9208c3e958416444ea23922ce5d260c838da2437dadb5e5f110d6877}
\end{figure}

\begin{figure}[htbp]
\centering
\includegraphics[width=\linewidth]{38203cc950fa0189b22b31bf9677bf1ed60c9fb522cd4ccd4b3fc68fb3088bc4.jpg}
\label{fig:38203cc950fa0189b22b31bf9677bf1ed60c9fb522cd4ccd4b3fc68fb3088bc4}
\end{figure}

(g)

(f)
(e)
\begin{figure}[htbp]
\centering
\includegraphics[width=\linewidth]{78e6ddafe7c9a558f683a314e4fc185ceb9b843eaef1d16d65526526611b7f65.jpg}
\label{fig:78e6ddafe7c9a558f683a314e4fc185ceb9b843eaef1d16d65526526611b7f65}
\end{figure}

(h)

\begin{figure}[htbp]
\centering
\includegraphics[width=\linewidth]{060b8e619f9c08bb2f70d19018afb69e81988d5c7556b91ecd2021e099208f4b.jpg}
\label{fig:060b8e619f9c08bb2f70d19018afb69e81988d5c7556b91ecd2021e099208f4b}
\end{figure}

(i)
图 3.9 使用鲁棒损失函数和渐进非凸性（GNC）求解含离群值的 SLAM 问题：(a)–(c) 为 M3500、SubT 和 Victoria Park 数据集的真值轨迹；(d)–(f) 为使用 Geman–McClure 损失得到的轨迹估计；(g)–(i) 为使用截断二次损失得到的轨迹估计。测量以彩色边显示：灰色边为正确分类的边；红色边为被错误归类为内点的离群值（假阳性）；蓝色边为被错误归类为离群值的内点（假阴性）。

M3500
\begin{figure}[htbp]
\centering
\includegraphics[width=\linewidth]{d2882ed1424ef18300c55e7d00e57729eb0b84a2385d828ac2329d94ed020750.jpg}
\label{fig:d2882ed1424ef18300c55e7d00e57729eb0b84a2385d828ac2329d94ed020750}
\end{figure}

(a)

SubT
\begin{figure}[htbp]
\centering
\includegraphics[width=\linewidth]{5df8cba2e5315e1f2caec4c78b5e47f41ddf1b2c5205778b831bd48f8663c940.jpg}
\label{fig:5df8cba2e5315e1f2caec4c78b5e47f41ddf1b2c5205778b831bd48f8663c940}
\end{figure}

Victoria Park

\begin{figure}[htbp]
\centering
\includegraphics[width=\linewidth]{e09061bf60bd4e86e3bbc576d7dd23edebb5227a069c30cd3eec5248284cba8d.jpg}
\label{fig:e09061bf60bd4e86e3bbc576d7dd23edebb5227a069c30cd3eec5248284cba8d}
\end{figure}

\begin{figure}[htbp]
\centering
\includegraphics[width=\linewidth]{fa5804728b14230829381ece9f392d23b9b0ec5990742ae2c6676d1649039eb1.jpg}
\label{fig:fa5804728b14230829381ece9f392d23b9b0ec5990742ae2c6676d1649039eb1}
\end{figure}

(b)
\begin{figure}[htbp]
\centering
\includegraphics[width=\linewidth]{a7c09461b4a23d467b3bcb78ca6b10fa61d1ea1a63ea1b27814eee4b680fd151.jpg}
\label{fig:a7c09461b4a23d467b3bcb78ca6b10fa61d1ea1a63ea1b27814eee4b680fd151}
\end{figure}

(c)
\begin{figure}[htbp]
\centering
\includegraphics[width=\linewidth]{4445cccec0d0f23f181f9b78da1f6e410799d47cb933104b566421957008e984.jpg}
\label{fig:4445cccec0d0f23f181f9b78da1f6e410799d47cb933104b566421957008e984}
\end{figure}

(d)
\begin{figure}[htbp]
\centering
\includegraphics[width=\linewidth]{90e841f1a16ac212ff2be7066b957819a997ccc70de5f3501f3fd36f8ab00ec2.jpg}
\label{fig:90e841f1a16ac212ff2be7066b957819a997ccc70de5f3501f3fd36f8ab00ec2}
\end{figure}

(g)

(e)
\begin{figure}[htbp]
\centering
\includegraphics[width=\linewidth]{3c2bf85279f9d3bf422848becd821a386a8abe8ac06de369855d606430bde9d6.jpg}
\label{fig:3c2bf85279f9d3bf422848becd821a386a8abe8ac06de369855d606430bde9d6}
\end{figure}

(h)

(f)
\begin{figure}[htbp]
\centering
\includegraphics[width=\linewidth]{aaca4bf1d9cb5663a25d86c161fac9f126bdffe2719cb8e688575f342d59937f.jpg}
\label{fig:aaca4bf1d9cb5663a25d86c161fac9f126bdffe2719cb8e688575f342d59937f}
\end{figure}

(i)
图 3.10 回环、路标测量和里程计中均含离群值的 SLAM 问题：(a)–(c) 为使用截断二次损失的 GNC 得到的轨迹估计；(d)–(f) 为使用 PCM 进行前端离群值剔除、随后采用最小二乘优化得到的轨迹估计；(g)–(i) 为先使用 PCM 进行前端离群值剔除、再使用带截断二次损失的 GNC 得到的轨迹估计。

在里程计与回环测量均可能为离群值时，鲁棒地求解 SLAM 问题极其困难。为展示所讨论方法在这个更复杂设置下的不足，我们先通过随机扰动两条里程计测量，修改图 3.1 中的位姿图问题，然后尝试用 GNC 求解。具体地，在 GNC 中，除了两条可能被扰动的测量之外，所有里程计测量均固定为内点。图 3.10 展示了 GNC 求解得到的轨迹。由于初始化受到扰动，GNC 失败并收敛到一个局部极小值，将所有测量（包括内点）都归类为离群值（蓝色边）。同一图还展示了在采用最小二乘后端之前使用 PCM 过滤严重离群值的结果。PCM 与初始猜测无关，因此在 SubT 和 Victoria Park 数据集上能够收敛到更好的解（图 3.10(e) 和 (f)）。有趣的是，在 SubT 和 Victoria Park 数据集上，组合前端离群值剔除（PCM）与 GNC 获得了最佳结果。与此同时，三种方法（GNC、PCM 以及 PCM+GNC）在 M3500 数据集上均未能收敛到可接受的解，这证实了这一 SLAM 设置的困难性，以及现有 SLAM 算法在离群值剔除方面的局限。

\section{延伸阅读与近期趋势}

\textbf{一致性最大化。} 尽管本章讨论的是 RANSAC 算法最基本的实例（依照 \textsuperscript{[328]} 的最初提议），但值得一提的是，文献中给出了许多 RANSAC 变体，包括通过局部优化精化估计的变体 \textsuperscript{[216]}、在 RANSAC 迭代中使用更好评分（而非一致集大小）的变体（例如 MLESAC \textsuperscript{[872]}），以及在 RANSAC 迭代中对采样施加偏置的变体（例如 PROSAC \textsuperscript{[215]}）。近期文献还包含可微 RANSAC 变体 \textsuperscript{[1174]}，以及在式（3.4）中的参数 $\gamma$ 未知时尝试寻找内点的变体（例如 \textsuperscript{[46]}）。文献 \textsuperscript{[1109]} 对 RANSAC 变体进行了近期综述与评估。

除 RANSAC 外，文献还包含精确一致性最大化方法，通常基于分支定界 \textsuperscript{[64, 436, 1290, 651, 1026, 203, 507, 132, 1226, 1224]}。尽管具有全局最优性保证，分支定界在最坏情形下具有指数运行时间，无法扩展到高维问题。

\textbf{成对一致性最大化。} 第 3.2.2 节所述的 PCM 方法最初由 Mangelson 等人 \textsuperscript{[728]} 在多机器人 SLAM 的语境下提出，并显示出特别的潜力。计算机视觉中也研究了基于图论的离群值剔除。San Segundo 和 Artieda \textsuperscript{[967]} 构建关联图，并为二维图像特征匹配寻找最大团。Perera 和 Barnes \textsuperscript{[869]} 用团形式化方法分割刚体运动下的对象。Leordeanu 和 Hebert \textsuperscript{[639]} 通过近似谱方法，在对应图中寻找强连接簇来建立图像匹配。Enqvist 等人 \textsuperscript{[309]} 开发了基于近似点覆盖的三维–三维和二维–三维配准离群值剔除算法。Yang 和 Carlone \textsuperscript{[1217]}、Yang 等人 \textsuperscript{[1223]} 与 Bustos 等人 \textsuperscript{[133]} 研究了基于最大团的三维–三维配准图论离群值剔除。检查一个测量子集内部一致性的思想也出现在 Latif 等人 \textsuperscript{[628]} 中：该工作通过聚类测量，并用基于卡方的检验检查一致性，从而剔除回环离群值。PCM 论文 \textsuperscript{[728]} 与本章讨论一样聚焦于成对一致性。近期，PCM 已在 \textsuperscript{[1005, 1006, 332]} 中推广到群组 $k$ 一致性，即一致性约束（3.7）涉及 $k$ 个测量而非仅两个测量的情形。这些论文本质上将一致性图的概念推广为一致性超图，其中每条超边涉及 $k$ 个节点。相关工作还考虑最大团问题的软变体：将二元条件（3.7）放宽，以在一致性图的边上产生连续权重 \textsuperscript{[712, 711]}。这些方法已应用于包括地下探索 \textsuperscript{[298]}、LiDAR 点云定位 \textsuperscript{[713]}、多机器人度量–语义建图 \textsuperscript{[1095]} 和非结构化环境全局定位 \textsuperscript{[33]} 在内的实际应用。

\textbf{交替最小化与渐进非凸性。} M-估计一直是机器人学 \textsuperscript{[99]} 和视觉 \textsuperscript{[978, 173]} 中常用的鲁棒估计方法。Tavish 和 Barfoot \textsuperscript{[1075]} 研究了不同损失函数的性能。若干论文研究了式（3.16）那样带辅助变量的形式化，却未认识到 Black–Rangarajan 对偶性（定理 3.4）所提供的与 M-估计的联系。例如，Sunderhauf 和 Protzel \textsuperscript{[1059]}、Sunderhauf 和 Protzel \textsuperscript{[1058]} 以及 Agarwal 等人 \textsuperscript{[11]} 用负责禁用离群值的潜在二元变量扩展问题。Lee 等人 \textsuperscript{[637]} 使用期望最大化。Olson 和 Agarwal \textsuperscript{[827]} 使用最大混合分布近似多模态测量噪声。近期，Barron \textsuperscript{[57]} 提出一个单参数函数，概括 M-估计中一族鲁棒损失函数。Chebrolu 等人 \textsuperscript{[174]} 设计了一个期望最大化算法，以同时估计未知量 $x$ 并选择合适的鲁棒损失函数 $\rho$。渐进非凸性算法最初由 \textsuperscript{[86, 85]} 为早期视觉应用中的离群值剔除提出；近期该算法已用于点云配准 \textsuperscript{[1297, 1223]}、SLAM \textsuperscript{[1222]} 和其他应用 \textsuperscript{[34]}。近期，Peng 等人 \textsuperscript{[860]} 提出了一种与 GNC 和 IRLS 相似的算法，基于优化中的平滑主化（smooth majorization）思想，适用于广泛的鲁棒损失。此外，Peng 等人 \textsuperscript{[860]} 还为 GNC 推导了全局和局部收敛保证。

\textbf{可认证算法。} 至此所述算法大致可分为两类：(i) 快速启发式方法（例如 RANSAC 或 M-估计的局部求解器），它们高效但性能保证很少；(ii) 全局求解器（例如分支定界），它们提供最优性保证但难以随问题规模扩展。近年来出现了一类新方法，称为\textbf{可认证算法}（certifiable algorithms），试图在可处理性与最优性之间取得平衡。可认证算法将非凸鲁棒估计问题松弛为凸半定规划（semidefinite programs，SDP），\textsuperscript{15} 其解可在多项式时间内获得，并提供可直接检验的事后全局最优性证书。鲁棒估计的可认证算法已在旋转估计 \textsuperscript{[1218]}、三维–三维配准 \textsuperscript{[1223]} 和位姿图优化 \textsuperscript{[618, 153]} 的语境下提出。文献 \textsuperscript{[1219, 1221]} 描述了一种为含离群值问题推导可认证算法的相当通用的方法，而 \textsuperscript{[147]} 讨论了其与统计学并行工作的联系。除少数显著例外外，这些算法尽管在多项式时间内运行，计算代价仍高，且通常比启发式方法慢得多。在一些情形中，这些算法背后的洞见可用来认证快速启发式方法获得的解的最优性 \textsuperscript{[1219]}，从而兼得二者之长。
